\documentclass[aspectratio=169, 10pt]{beamer}
\usepackage[utf8]{inputenc}
\usepackage[spanish]{babel}
\usepackage{tikz}
\usepackage{booktabs}
\usepackage{xcolor}
\usepackage{fontawesome5}
\usetikzlibrary{shapes.geometric, arrows.meta, positioning, fit, backgrounds}

\usetheme{Madrid}
\usecolortheme{default}
\setbeamercolor{palette primary}{bg=blue!60!black, fg=white}
\setbeamercolor{palette secondary}{bg=blue!40!black, fg=white}
\setbeamercolor{palette tertiary}{bg=blue!20!black, fg=white}
\setbeamercolor{structure}{fg=blue!60!black}
\setbeamercolor{block title}{bg=blue!15, fg=blue!60!black}
\setbeamercolor{alerted text}{fg=orange!80!black}

\title[Modelo Relacional SISBEN]{Arquitectura del Modelo Sem\'antico\\Tablero SISBEN -- Antioquia}
\subtitle{Dise\~no relacional, tablas de par\'ametro y organizaci\'on de medidas}
\author{Departamento de Planeaci\'on de Antioquia}
\date{Mayo 2026}

\newcommand{\yes}{\textcolor{green!60!black}{\faCheck}}
\newcommand{\no}{\textcolor{red!70!black}{\faTimes}}
\newcommand{\warn}{\textcolor{orange}{\faExclamationTriangle}}

\begin{document}

\begin{frame}
  \titlepage
\end{frame}

\begin{frame}{Contenido}
  \tableofcontents
\end{frame}

\section{Inventario del Modelo}

\begin{frame}{Inventario general del modelo}
\begin{columns}
\column{0.48\textwidth}
\begin{block}{Objetos del modelo (mayo 2026)}
\begin{tabular}{lr}
\toprule
Objeto & Cantidad \\
\midrule
Tablas totales & 18 \\
\quad Tablas de hechos & 2 \\
\quad Dimensiones conectadas & 4 \\
\quad Tablas de par\'ametro & 10 \\
Columnas (fct\_Sisben) & 135 \\
Medidas totales & 203 \\
Relaciones activas & 6 \\
Carpetas de medidas & 7 \\
\bottomrule
\end{tabular}
\end{block}

\column{0.48\textwidth}
\begin{block}{Distribuci\'on de medidas por carpeta}
\begin{tabular}{lr}
\toprule
Carpeta & Medidas \\
\midrule
IPM + Dim. IPM & $\sim$55 \\
Vivienda (3 subcarpetas) & $\sim$50 \\
Hogares + Gasto & $\sim$35 \\
Demograf\'ia & $\sim$30 \\
Selectores & 12 \\
Calidad Datos & 13 \\
UX & 6 \\
\bottomrule
\end{tabular}
\end{block}
\end{columns}
\end{frame}

\section{Esquema Relacional}

\begin{frame}{Esquema estrella del modelo}
\centering
\begin{tikzpicture}[
  fact/.style={rectangle, draw=blue!60!black, fill=blue!15, rounded corners=3pt,
               minimum width=2.8cm, minimum height=0.9cm, font=\small\bfseries},
  dim/.style={rectangle, draw=green!50!black, fill=green!10, rounded corners=3pt,
              minimum width=2.5cm, minimum height=0.7cm, font=\scriptsize},
  arr/.style={-{Stealth[length=5pt]}, thick, gray!70},
  node distance=0.4cm
]

% Tablas de hechos (centro)
\node[fact] (sisben) {fct\_Sisben\\{\tiny 135 cols | 191 medidas}};
\node[fact, below=0.6cm of sisben] (pob) {Poblacion\_DANE\\{\tiny 6 cols | 1 medida}};

% Dimensiones conectadas
\node[dim, right=2.5cm of sisben, yshift=0.5cm] (mpio) {dim\_municipios\\{\tiny 21 cols}};
\node[dim, right=2.5cm of sisben, yshift=-0.5cm] (est) {Dim\_Estado\_DNP\\{\tiny 3 cols}};
\node[dim, left=2.5cm of sisben] (anios) {Dim\_Anios\\{\tiny 1 col}};
\node[dim, right=2.5cm of pob, yshift=-0.3cm] (rango) {Dim\_Rango\_Edad\\{\tiny 2 cols}};

% Relaciones
\draw[arr] (sisben) -- node[above, font=\tiny]{N:1} (mpio);
\draw[arr] (sisben) -- node[below, font=\tiny]{N:1} (est);
\draw[arr] (sisben) -- node[above, font=\tiny]{N:1} (anios);
\draw[arr] (pob) -- node[above, font=\tiny]{N:1} (mpio);
\draw[arr] (pob) -- node[below, font=\tiny]{N:1} (rango);
\draw[arr] (pob) to[bend left=20] node[left, font=\tiny]{N:1} (anios);

\end{tikzpicture}

\vspace{0.3cm}
{\small \alert{Esquema estrella cl\'asico}: cada tabla de hechos conectada a dimensiones via claves. Todas las relaciones son \textbf{N:1 activas, filtraje en una direcci\'on}.}
\end{frame}

\section{Tablas de Par\'ametro Desconectadas}

\begin{frame}{?`Qu\'e son las tablas de par\'ametro?}
\begin{columns}
\column{0.52\textwidth}
\textbf{Definici\'on:} tablas sin relaci\'on con la tabla de hechos, usadas exclusivamente como \textit{slicer} para controlar qu\'e medida se muestra en un visual.

\vspace{0.4cm}
\begin{alertblock}{Patr\'on DAX: Disconnected Slicer Table}
\begin{enumerate}\footnotesize
  \item La tabla tiene 1 columna con nombres de opciones
  \item El usuario selecciona una opci\'on en el slicer
  \item La medida captura esa selecci\'on con \texttt{SELECTEDVALUE()}
  \item Un \texttt{SWITCH()} devuelve el c\'alculo correspondiente
\end{enumerate}
\end{alertblock}

\column{0.44\textwidth}
\begin{block}{Tablas de par\'ametro en el modelo}
\begin{tabular}{ll}
\toprule
Tabla & Controla \\
\midrule
Nivel\_Sisben & Eje de la matriz \\
Dim\_Privacion & Tipo de privaci\'on \\
Dim\_IPM\_Dimension & Indicador IPM \\
Dim\_Evento & Evento natural \\
Dim\_Indicador\_Hogar & Indicador hogar \\
Dim\_Bloque\_Sanitario & Variable sanitaria \\
Dim\_Categoria\_Vivienda & Categor\'ia material \\
Selector\_Vivienda & Tipo de vivienda \\
Tabla\_Ind\_Vivienda & Servicio p\'ublico \\
\bottomrule
\end{tabular}
\end{block}
\end{columns}
\end{frame}

\begin{frame}[fragile]{C\'omo funciona: ejemplo real del modelo}

\begin{columns}
\column{0.48\textwidth}
\textbf{Tabla \texttt{Dim\_IPM\_Dimension}} (slicer):
\begin{center}
\begin{tikzpicture}
\node[draw=gray, fill=gray!10, rounded corners, font=\scriptsize,
      inner sep=6pt, align=left]{
  Dimension\_IPM \\ \midrule
  Bajo logro educativo \\
  Analfabetismo \\
  Inasistencia escolar \\
  Rezago escolar \\
  Barreras primera infancia \\
  \ldots (15 indicadores)
};
\end{tikzpicture}
\end{center}

\column{0.50\textwidth}
\textbf{Medida \texttt{Pct\_Privacion\_Seleccionada}:}
\vspace{0.2cm}
\begin{tikzpicture}
\node[draw=blue!30, fill=blue!5, rounded corners, font=\tiny\ttfamily,
      inner sep=6pt, align=left, text width=5.8cm]{
  \textbf{VAR} Val = SWITCH(\newline
  \quad SELECTEDVALUE(\newline
  \quad\quad Dim\_IPM\_Dimension[Dimension\_IPM]),\newline
  \quad "Bajo logro educativo",\newline
  \quad\quad [Pct\_I1\_BajoLogroEducativo],\newline
  \quad "Analfabetismo",\newline
  \quad\quad [Pct\_I2\_Analfabetismo],\newline
  \quad ...\newline
  )\newline
  \textbf{RETURN} Val
};
\end{tikzpicture}
\end{columns}

\vspace{0.3cm}
\centering\footnotesize
\textbf{Resultado:} una sola columna en la matriz muestra cualquiera de los 15 indicadores IPM seg\'un la selecci\'on del usuario.
\end{frame}

\begin{frame}{Flujo completo: del slicer al visual}
\centering
\begin{tikzpicture}[
  box/.style={rectangle, rounded corners=4pt, draw, fill=white,
              minimum height=0.7cm, minimum width=2.8cm, font=\small, align=center},
  slicer/.style={box, fill=orange!15, draw=orange!60!black},
  measure/.style={box, fill=blue!12, draw=blue!50!black},
  visual/.style={box, fill=green!12, draw=green!50!black},
  arr/.style={-{Stealth[length=5pt]}, thick},
  note/.style={font=\scriptsize, gray}
]

\node[slicer] (s1) at (0,0) {Slicer\\{\tiny Dim\_IPM\_Dimension}};
\node[note] at (0,-0.7) {usuario elige};
\node[note] at (0,-1.0) {"Analfabetismo"};

\node[measure] (m1) at (4.5,0) {\texttt{SELECTEDVALUE()}\\{\tiny captura selecci\'on}};

\node[measure] (m2) at (8.5,0) {\texttt{SWITCH()}\\{\tiny mapea a medida base}};

\node[measure] (m3) at (12,0) {\texttt{[Pct\_I2]}\\{\tiny calcula sobre fct\_Sisben}};

\node[visual] (v1) at (6.5,-2) {Matriz / Visual\\{\tiny muestra el valor}};

\draw[arr] (s1) -- (m1);
\draw[arr] (m1) -- (m2);
\draw[arr] (m2) -- (m3);
\draw[arr, bend right=20] (m3) to (v1);
\draw[arr, bend left=20] (m1) to (v1);

\draw[dashed, gray!60] (6.5, 0.8) -- (6.5,-1.3) node[right, font=\scriptsize, gray]{flujo de evaluaci\'on};

\end{tikzpicture}

\vspace{0.3cm}
{\footnotesize La tabla de par\'ametro \textbf{nunca se une a fct\_Sisben}. Solo proporciona el contexto de filtro que la medida \texttt{SELECTEDVALUE()} captura en cada evaluaci\'on.}
\end{frame}

\begin{frame}{?`Por qu\'e no hacer columnas directamente en la matriz?}

\begin{columns}
\column{0.48\textwidth}
\begin{exampleblock}{Sin tablas de par\'ametro (opci\'on alternativa)}
\begin{itemize}\footnotesize
  \item Se agregar\'ian directamente las columnas \texttt{I1}, \texttt{I2}\ldots \texttt{I15} a la matriz
  \item El usuario no puede cambiar qu\'e columnas ve
  \item La matriz se vuelve \textbf{horizontal e ilegible} (15 columnas)
  \item No es posible comparar un indicador vs. municipios en filas
  \item Imposible construir visuals din\'amicos con un solo KPI cambiante
\end{itemize}
\end{exampleblock}

\column{0.48\textwidth}
\begin{alertblock}{Con tablas de par\'ametro (opci\'on implementada)}
\begin{itemize}\footnotesize
  \item[\yes] La matriz siempre tiene \textbf{una sola columna de valor}
  \item[\yes] El usuario elige qu\'e indicador ver con el slicer
  \item[\yes] Compatible con cualquier combinaci\'on de filtros geogr\'aficos
  \item[\yes] Permite comparar municipios (filas) vs. un indicador a la vez
  \item[\yes] Es el patr\'on est\'andar recomendado por Microsoft para matrices din\'amicas
\end{itemize}
\end{alertblock}
\end{columns}
\end{frame}

\section{Organizaci\'on de Medidas}

\begin{frame}{Jerarqu\'ia de carpetas (Display Folders)}
\begin{columns}
\column{0.50\textwidth}
\begin{tikzpicture}[
  folder/.style={draw=gray!60, fill=gray!8, rounded corners=2pt,
                 font=\scriptsize, inner sep=4pt, minimum width=4.8cm, align=left},
  sub/.style={draw=blue!30, fill=blue!5, rounded corners=2pt,
              font=\scriptsize, inner sep=3pt, minimum width=4.3cm, align=left},
  node distance=0.18cm
]
\node[folder](f1){IPM \quad{\tiny ($\sim$40 medidas)}};
\node[sub, below=0.1cm of f1, xshift=0.3cm](f1a){Dimensiones IPM \quad{\tiny (15)}};
\node[folder, below=0.22cm of f1a, xshift=-0.3cm](f2){Vivienda \quad{\tiny ($\sim$30)}};
\node[sub, below=0.1cm of f2, xshift=0.3cm](f2a){Privaciones \quad{\tiny (10)}};
\node[sub, below=0.1cm of f2a](f2b){Eventos Naturales \quad{\tiny (11)}};
\node[folder, below=0.22cm of f2b, xshift=-0.3cm](f3){Hogares \quad{\tiny (16)}};
\node[sub, below=0.1cm of f3, xshift=0.3cm](f3a){Gasto \quad{\tiny (17)}};
\node[folder, below=0.22cm of f3a, xshift=-0.3cm](f4){Demograf\'ia \quad{\tiny ($\sim$30)}};
\node[folder, below=0.22cm of f4](f5){Selectores \quad{\tiny (12)}};
\node[folder, below=0.22cm of f5](f6){Calidad Datos \quad{\tiny (13)}};
\node[folder, below=0.22cm of f6](f7){UX \quad{\tiny (6)}};
\end{tikzpicture}

\column{0.46\textwidth}
\begin{block}{Criterio de clasificaci\'on}
\begin{description}\footnotesize
  \item[\textbf{IPM}] \'Indices, incidencias, rankings, variaciones, textos
  \item[\textbf{Vivienda}] Condiciones f\'isicas, privaciones, riesgo
  \item[\textbf{Hogares}] Bienes, servicios, gasto del hogar
  \item[\textbf{Demograf\'ia}] Conteos, grupos SISBEN, pir\'amides
  \item[\textbf{Selectores}] Medidas que dependen de una tabla de par\'ametro
  \item[\textbf{Calidad}] Validaci\'on, matrices de control
  \item[\textbf{UX}] T\'itulos din\'amicos, colores, etiquetas
\end{description}
\end{block}

\vspace{0.2cm}
\footnotesize Las medidas de las tablas de par\'ametro se clasifican en \textbf{Selectores}, indicando expl\'icitamente que su valor depende de un slicer externo.
\end{columns}
\end{frame}

\section{Resumen y Justificaci\'on}

\begin{frame}{Resumen: decisiones de dise\~no justificadas}

\begin{block}{Principios aplicados}
\begin{tabular}{p{5.5cm} p{6cm}}
\toprule
\textbf{Decisi\'on} & \textbf{Justificaci\'on} \\
\midrule
Esquema estrella con 6 relaciones N:1 & M\'inimo de relaciones necesarias; rendimiento \'optimo en VertiPaq \\
\addlinespace
10 tablas de par\'ametro desconectadas & Patr\'on est\'andar para matrices din\'amicas; evita proliferaci\'on de visuals \\
\addlinespace
\texttt{SWITCH(SELECTEDVALUE())} como motor & Evaluaci\'on diferida; solo se computa el c\'alculo seleccionado \\
\addlinespace
203 medidas en 7 carpetas jer\'arquicas & Mantenibilidad; navegaci\'on clara en Power BI Desktop \\
\addlinespace
\texttt{fct\_Sisben} como nombre can\'onico & Convenci\'on \texttt{fct\_/dim\_} est\'andar de modelado tabular \\
\bottomrule
\end{tabular}
\end{block}

\begin{exampleblock}{Conclusi\'on}
Las tablas aisladas \textbf{no son un error de dise\~no}: son el mecanismo que hace posible que cada visual del tablero sea interactivo, flexible y din\'amico, sin duplicar c\'odigo DAX ni multiplicar el n\'umero de visuals necesarios.
\end{exampleblock}
\end{frame}

\begin{frame}{Mapa completo del modelo}
\centering
\begin{tikzpicture}[
  fact/.style={rectangle, draw=blue!60!black, fill=blue!15, rounded corners=3pt,
               minimum width=2.5cm, minimum height=0.75cm, font=\scriptsize\bfseries},
  dim/.style={rectangle, draw=green!50!black, fill=green!10, rounded corners=3pt,
              minimum width=2.2cm, minimum height=0.6cm, font=\tiny},
  slicer/.style={rectangle, draw=orange!60!black, fill=orange!10, rounded corners=3pt,
                 minimum width=2.2cm, minimum height=0.55cm, font=\tiny},
  arr/.style={-{Stealth[length=4pt]}, thick, gray!70},
  node distance=0.25cm
]
% Fact tables
\node[fact] (fct) at (0,0) {fct\_Sisben};
\node[fact] (pop) at (0,-2.2) {Poblacion\_DANE};

% Connected dims
\node[dim] (mpio) at (3.8, 0.8) {dim\_municipios};
\node[dim] (est) at (3.8, 0.0) {Dim\_Estado\_DNP};
\node[dim] (anio) at (3.8, -0.8) {Dim\_Anios};
\node[dim] (rango) at (3.8, -2.2) {Dim\_Rango\_Edad};

% Arrows
\draw[arr] (fct) -- (mpio);
\draw[arr] (fct) -- (est);
\draw[arr] (fct) -- (anio);
\draw[arr] (pop) -- (mpio);
\draw[arr] (pop) -- (anio);
\draw[arr] (pop) -- (rango);

% Slicer tables
\node[slicer] (s1) at (-4.2, 1.8) {Nivel\_Sisben};
\node[slicer] (s2) at (-4.2, 1.1) {Dim\_Privacion};
\node[slicer] (s3) at (-4.2, 0.4) {Dim\_IPM\_Dimension};
\node[slicer] (s4) at (-4.2,-0.3) {Dim\_Evento};
\node[slicer] (s5) at (-4.2,-1.0) {Dim\_Indicador\_Hogar};
\node[slicer] (s6) at (-4.2,-1.7) {Dim\_Bloque\_Sanitario};
\node[slicer] (s7) at (-4.2,-2.4) {Dim\_Categoria\_Vivienda};
\node[slicer] (s8) at (-4.2,-3.1) {Selector\_Vivienda};
\node[slicer] (s9) at (-4.2,-3.8) {Tabla\_Ind\_Vivienda};

% Dashed connection label
\draw[dashed, orange!50, very thick] (-2.8,2.2) -- (-2.8,-4.1);
\node[rotate=90, font=\tiny, orange!70!black] at (-2.5,-1.0) {sin relaci\'on -- solo slicer};

% Legend
\node[fact, minimum width=1.5cm, minimum height=0.45cm, font=\tiny] at (1.2,-4.3) {};
\node[font=\tiny, right] at (2.1,-4.3) {Tabla de hechos};
\node[dim, minimum width=1.5cm, minimum height=0.45cm, font=\tiny] at (4.0,-4.3) {};
\node[font=\tiny, right] at (4.9,-4.3) {Dimensi\'on};
\node[slicer, minimum width=1.5cm, minimum height=0.45cm, font=\tiny] at (6.8,-4.3) {};
\node[font=\tiny, right] at (7.7,-4.3) {Tabla de par\'ametro};

\end{tikzpicture}
\end{frame}

\end{document}
